% Generated draft: requires mathematical and source review.
\documentclass[11pt]{article}
\usepackage[margin=25mm]{geometry}
\usepackage{fontspec}
\setmainfont{Latin Modern Roman}
\usepackage{amsmath,amssymb,mathtools,unicode-math}
\setmathfont{Latin Modern Math}
\usepackage{enumitem,needspace,xurl,hyperref}
\setlength{\parindent}{0pt}
\setlength{\parskip}{5pt}
\setlength{\emergencystretch}{3em}
\newcommand{\N}{\mathbb{N}}
\newcommand{\Z}{\mathbb{Z}}
\newcommand{\Q}{\mathbb{Q}}
\newcommand{\R}{\mathbb{R}}
\newcommand{\C}{\mathbb{C}}
\newcommand{\F}{\mathbb{F}}
\newcommand{\A}{\mathbb{A}}
\newcommand{\PP}{\mathbb{P}}
\newcommand{\E}{\mathbb{E}}
\newcommand{\eps}{\varepsilon}
\newcommand{\dd}{\,\mathrm{d}}
\begin{document}
\section*{UnsolvedMath ID: 30003457}
\textbf{Source problem code:} catalogue entry
\subsection{Definitions and mathematical statement}
The supplied record concerns a rank-\(n\) vector bundle \(E\) on \(S\) and its classifying map \(f_E:S\to BGL_n\). It does not specify the base, the category containing \(S\), the topology used for the classifying stack, or the map \(\mathrm{pt}\to BGL_n\). Consequently, the displayed fiber product requires additional conventions.

One precise conditional formulation is the following. Fix a base scheme \(B\), a \(B\)-scheme \(S\), and a finite locally free \(\mathcal O_S\)-module \(E\) of constant rank \(n\). Interpret \(BGL_{n,B}\) as the fppf classifying stack of rank-\(n\) vector bundles, and interpret \(\mathrm{pt}\to BGL_n\) as the morphism \(B\to BGL_{n,B}\) classifying \(\mathcal O_B^{\oplus n}\). These are explicit supplementary conventions, not hypotheses specified in the source record.

Define the frame functor of \(E\) on \(S\)-schemes by
\[
\operatorname{Fr}_S(E)(T\xrightarrow{g}S)
=
\operatorname{Isom}_{\mathcal O_T}
\bigl(\mathcal O_T^{\oplus n},g^*E\bigr).
\]
The recoverable clarification asks for the geometric interpretation of the stack fiber product, with the requested identification
\[
S\mathbin{\times^{(2)}_{BGL_{n,B}}}B
\simeq
\operatorname{Fr}_S(E),
\]
where \(\times^{(2)}\) denotes the \(2\)-fiber product and the frame functor is viewed as its associated fppf sheaf or stack over \(S\). This states the identification requested by the supplied record; no proof is provided.

The source qualifications prevent promoting this clarification to a verified open problem. The registry labels the item \texttt{open}, but its supplied statement-verification note says that the official report immediately identifies the fiber product with the frame bundle and treats the question as an answered clarification. Its statement status is \texttt{unrecoverable}, with no clean statement supplied. The dated literature assessment instead describes the extracted text as the meta-question ``How to make this geometric?'' in a discussion of stacks, says that it is not a complete mathematical claim, and assigns no definite resolution. These supplied descriptions are retained without independently adjudicating them.
\subsection{Short English statement}
Under explicitly stated scheme and stack conventions, the question asks how pulling back the trivial-bundle point along the classifying map of \(E\) describes the bundle of frames of \(E\). The supplied record contains conflicting status descriptions: an \texttt{open} label, an assessment of an incomplete meta-question, and a verification note describing an already answered clarification. This draft therefore records a conditional formulation requiring human review.
\subsection{Sources}
\begin{itemize}
\item Local UnsolvedMath registry: supplied complete record for permanent ID \(30003457\), titled \emph{Frame Bundles as Stack Fiber Products}; exact source problem code \texttt{catalogue entry}, in the Oberwolfach Reports Open Problems collection. The proposer is unspecified; the supplied proposed year is \(2017\).
\item Supplied report citation: \emph{Arbeitsgemeinschaft: Higher Gross Zagier Formulas (2018)}, Oberwolfach Reports, pp.~1067--1134, DOI \texttt{10.4171/owr/2017/18}. Supplied source URL and literature-sources link, labeled ``Original Oberwolfach report'': \url{https://doi.org/10.4171/owr/2017/18}. The background attributes the collection to workshop reports and their open-problem discussions.
\item Literature assessment checked on \texttt{2026-08-21}: this is supplied dated curation, not independently verified. It labels the item open while describing ``How to make this geometric?'' as an incomplete meta-question about stacks for which no definite resolution can be assigned. The background requests independent verification of time-sensitive status and cited claims.
\item The supplied statement-verification note attributes the frame-bundle identification directly to the official report and describes the extracted question as an answered clarification, not an open problem. This attribution has not been independently checked here.
\end{itemize}
\end{document}
